\section{模型评价与进一步讨论}\label{sec:evaluation}

\subsection{三个问题的衔接与适用范围}
本文按照“可追溯特征构建—缺失场景预测—决策证据解释”组织三项工作。问题一在附件1原始媒体上建立统一时间窗和来源链，保证特征位置有明确的媒体时间依据；问题二使用附件2已经提供的对齐特征，检验轻量预测模型在有效位置局部连续缺失时的开发验证表现；问题三将确定的预测函数用于附件4，生成模态贡献、交互、关键特征槽位与可核验文本片段。三项工作共享掩码和时间语义的审慎原则，但不应被误写为附件1特征直接输入问题二模型。

\subsection{预测收益、复杂度与误差边界}
问题二的改进结构仅比基线增加883个参数，三次随机初始化的平均综合鲁棒分数增加0.003100，但其中一次基本持平。平均缺失条件下的宏平均F1、MAE及相关系数方向较好，准确率略低；部分双模态缺失组合的改进接近零，原生视觉全零样本也存在不利结果。因此，较小的参数开销和平均分数收益是本方法的主要优势，而非对所有缺失类型的普遍优势。基线按完整输入分数、改进模型按综合分数选取参数的口径差异也限制了二者的严格横向解释。综合分数只是内部验证工具，正式判断仍须同时查看四项原始指标。

问题一的重复分层五折评价用于比较特征构建与对齐方案的信息保持能力。该评价的训练/评价折间存在同一视频标识的片段，因此不能将这些数值解释为独立视频上的泛化性能或统计显著性证据。问题二使用附件2预定训练与验证划分，所有模型选择及缺失场景比较限于验证集；附件2测试划分与附件3不支撑本文所报告的验证结论。

\subsection{解释证据的层级边界}
问题三的八联盟精确计算保证了模态贡献满足加和关系；连续遮挡和删除曲线提供预测函数内部的干预响应证据。然而关键窗口由最大响应准则选取，故其相对随机窗口的优势包含选择效应。文本特征行到标记、再到原始字符片段的链路在附件4得到数值核验，语音与视觉特征槽位仍无可靠的媒体时间映射。模型解释可以定位“哪种模态、哪个特征槽位影响当前输出”，但不能据此直接断言真实情绪因果、精确音频秒数或关键视频帧。附件4缺乏标签，也无法评估其预测准确性。

在实际应用中，可将类别与强度预测、模态贡献、局部敏感性及证据核验状态同时输出，使用户区分模型决策依据与原始媒体证据。后续若获取可验证的音视频预处理记录和原始时间映射，才能把当前特征槽位层的解释进一步扩展到音频时间和视频帧。
